\section{Domaine, grandeurs et portée logique du certificat computationnel des dérivations intrinsèques}

\begin{longtable}{@{}>{\raggedright\arraybackslash}p{0.25\textwidth}>{\raggedright\arraybackslash}p{0.65\textwidth}@{}}
\toprule
Bloc & Controles ejecutados\\
\midrule
\endhead
Matrice & 58 lignes, 15 champs, identifiants uniques et portée déclarée\\
Vac\'io & ordre de \(q_\pm\), positivité, valeurs propres, quatre caractères, identités constitutives et complétion du mode 30\\
Electrod\'ebil & \(D_A=q_+q_-\) y \(\sin^2\theta_W=1-D_A\)\\
Chaos & somme quadratique, \(899=30^2-1=29\cdot31\), coefficient quartique, profil fermé et résidus numériques\\
Gauss & Lévy, entropie et Lochs pour les bases 10 et 729\\
Lacunes & fenêtre exhaustive de \(10^8\) positions, lacunes \(21/22\), retours \(6/7\), intervalle certifié et empreintes du vérificateur multiprécision\\
Spectre & caractère \(\chi_{12}\), \(L(1)\), \(L(2)\), unité hyperbolique et tour de Bendersky\\
Hamiltonien modulaire & \(\Dmod\), poids 90/120, entropie, défaut orienté, inversion aire--mémoire, TFD et première loi\\
Type III & groupe des rapports \(30\Dmod\Z\) et \(\lambda_\partial=\exp(-30\Dmod)\) sous les hypothèses déclarées\\
CKM & inverse des coordonnées, projecteur impair, réflexion rationnelle, déterminant et entrelacement\\
Cosmolog\'ia & élimination de \(n\) et loi exacte de puissance six\\
\bottomrule
\end{longtable}

\section{Valeurs numériques certifiées des invariants constitutifs, spectraux et modulaires}

\begin{longtable}{@{}>{\raggedright\arraybackslash}p{0.38\textwidth}>{\raggedright\arraybackslash}p{0.52\textwidth}@{}}
\toprule
Objet & Valeur certifiée\\
\midrule
\endhead
\(\widehat\mu\) & \(0.999448350479326935089981555906075895\ldots\)\\
\(\widehat\varepsilon\) & \(0.999999257096636702760441615542323693\ldots\)\\
\(\widehat Z\) & \(0.999724508538937215455554346613611395\ldots\)\\
\(\widehat c\) & \(1.000276310486157526106386268929576268\ldots\)\\
\(D_A\) & \(0.775129119247156430188884096480202784\ldots\)\\
\(1-D_A\) & \(0.224870880752843569811115903519797216\ldots\)\\
\(L(1,\chi_{12})\) & \(0.760345996300946347531094254880405824\ldots\)\\
\(L(2,\chi_{12})\) & \(0.949703126294009395263498491745741516\ldots\)\\
\(\gamma_{\mathbb Q(\sqrt3)}\) & \(1.0539656082484825800\ldots\)\\
\(L\) de Lévy & \(1.186569110415625452821722975947237121\ldots\)\\
\(h_G\) & \(2.373138220831250905643445951894474241\ldots\)\\
\(I_{\rm or}\) & \(0.001669183233868940655631225912996411\ldots\) nats\\
\(\lambda_\partial\) & \(0.996669646468957378610829976885164731\ldots\)\\
\bottomrule
\end{longtable}

\section{Obstructions de typage et séparation des équivalences entre domaines certifiés}

Le certificat ne promeut pas les affirmations suivantes :

\begin{enumerate}
\item \(I_{\rm or}=\gammaH a_0\) ; les valeurs sont proches mais distinctes.
\item L’opérateur du vide est identique au renormalisateur de Feigenbaum.
\item La cofinalité des cylindres implique une conjugaison TPK--Gauss.
\item La clôture \(III_{\lambda_\partial}\) vaut sans persistance des deux canaux.
\item L’inversion de feuillet CKM est la conjugaison CP physique.
\item La réalisation \(\varphi\) et la chronologie cyclique de \(108\) états
représentent la même cosmologie.
\item La relation angulaire dans le schéma des masses physiques inclut à elle seule toutes les corrections électrofaibles.
\end{enumerate}

Ces contrôles négatifs n’affaiblissent pas le volume : ils fixent exactement ce qui a été démontré et quelle application supplémentaire en changerait le statut.
